% --- Archon physics formalization source begin ---
% archon:physics
% archon:covers PhyXMiniProblems/problem_phyx_mini_0910.lean
% archon:source-report reports/phyx_mini/problem_phyx_mini_0910.source.json

\chapter{Physics problem phyx\_mini\_0910}
\label{ch:PhyXMiniProblems_problem_phyx_mini_0910}

\paragraph{Problem source.}
\#\# Physical scenario

Two \textbackslash{}( 1.0\textbackslash{},\textbackslash{}mathrm\{g\} \textbackslash{}) balls are connected by a \textbackslash{}( 2.0\textbackslash{},\textbackslash{}mathrm\{cm\} \textbackslash{})-long insulating rod of negligible mass. One ball has a charge of \textbackslash{}( +10\textbackslash{},\textbackslash{}mathrm\{nC\} \textbackslash{}), the other a charge of \textbackslash{}( -10\textbackslash{},\textbackslash{}mathrm\{nC\} \textbackslash{}). The rod is held in a \textbackslash{}( 1.0\textbackslash{}times10\textasciicircum{}\{4\}\textbackslash{},\textbackslash{}mathrm\{N/C\} \textbackslash{}) uniform electric field at an angle of \textbackslash{}( 30\textasciicircum{}\{\textbackslash{}circ\} \textbackslash{}) with respect to the field, then released.

\#\# Question

What is its initial angular acceleration?

\#\# Answer choices

- A: A: 3rad/s\textasciicircum{}2

- B: B: 2rad/s\textasciicircum{}2

- C: C: 5rad/s\textasciicircum{}2

- D: D: 6rad/s\textasciicircum{}2

\#\# Figure caption (auxiliary; use the image as primary evidence)

The image depicts a diagram involving two charged spheres connected by a rod. Here's a detailed breakdown:

1. **Spheres**: 
   - The left sphere is blue and negatively charged (\textbackslash{}( -10 \textbackslash{}, \textbackslash{}text\{nC\} \textbackslash{})) with a mass of \textbackslash{}( 1.0 \textbackslash{}, \textbackslash{}text\{g\} \textbackslash{}).
   - The right sphere is red and positively charged (\textbackslash{}( +10 \textbackslash{}, \textbackslash{}text\{nC\} \textbackslash{})) also with a mass of \textbackslash{}( 1.0 \textbackslash{}, \textbackslash{}text\{g\} \textbackslash{}).

2. **Rod**:
   - A thick rod connects the two spheres. The rod is shown at an angle of \textbackslash{}( 30\textasciicircum{}\textbackslash{}circ \textbackslash{}) relative to the horizontal line passing through the negative sphere.

3. **Measurements**:
   - The distance \textbackslash{}( s \textbackslash{}) between the centers of the two spheres is \textbackslash{}( 2.0 \textbackslash{}, \textbackslash{}text\{cm\} \textbackslash{}).
   - An electric field \textbackslash{}( E = 1.0 \textbackslash{}times 10\textasciicircum{}4 \textbackslash{}, \textbackslash{}text\{N/C\} \textbackslash{}) is depicted with horizontal pink arrows, pointing from left to right.
   - A dipole moment \textbackslash{}( \textbackslash{}vec\{p\} \textbackslash{}) is indicated along the direction of the rod, pointing from the negative to the positive charge.

4. **Arrows and Labels**:
   - Arrows indicate the directions of the electric field and the dipole moment.
   - Labels next to the spheres and measurements provide information about charge, mass, distance, and field strength.

This illustration likely represents a physics problem related to electric dipoles in an electric field.

\#\# Dataset metadata

Electrostatics; Physical Model Grounding Reasoning, Multi-Formula Reasoning

\paragraph{Recorded answer/context.}
C: C: 5rad/s\textasciicircum{}2

\paragraph{Figure/image path.}
/root/proposal\_for\_physic/science-mango/phyx\_mini\_run/phyx\_data/test\_image/910.png

\paragraph{Formalization target.}
create a compiling Lean file with sorry bodies at `PhyXMiniProblems/problem\_phyx\_mini\_0910.lean`. The Lean declarations must preserve the physical quantities, dimensions or dimensional roles, figure labels, governing-law hypotheses, and final relation expressed by this problem.
Use Mathlib/Physlib names found through LeanExplore where available. If a physics API is missing, introduce faithful local abstractions rather than scalar placeholder aliases.

\begin{theorem}[Physics formalization target]
\label{thm:physics:phyx_mini_0910:target}
\lean{PhyXMiniProblems.ProblemPhyXMini0910.initialAngularAcceleration_eq_fiveRadiansPerSecondSquared}
\uses{def:physics:phyx-mini-0910:phyxminiproblems-problemphyxmini0910-angularaccelerationinradianspersecondsquared, def:physics:phyx-mini-0910:phyxminiproblems-problemphyxmini0910-electricdipolereleasesetup, def:physics:phyx-mini-0910:phyxminiproblems-problemphyxmini0910-matchesprimaryfigure, def:physics:phyx-mini-0910:phyxminiproblems-problemphyxmini0910-matchesproblemdata, def:physics:phyx-mini-0910:phyxminiproblems-problemphyxmini0910-satisfieselectricdipolerotationallaws, def:physics:phyx-mini-0910:phyxminiproblems-problemphyxmini0910-recordeddatasetanswer, def:physics:phyx-mini-0910:phyxminiproblems-problemphyxmini0910-answermatchesinitialangularacceleration}
The assigned autoformalize agent should translate this physics problem into Lean declarations in the covered file, with theorem and lemma proof bodies written as `by sorry`.
\end{theorem}
\begin{proof}
This is an autoformalization task, not a proof task. Produce statements that can later be proved by the physics prover without weakening the physical model.
\end{proof}

% --- Archon declaration topology begin ---
\section{Declaration topology}
\begin{definition}[electric Field Strength Dimension]
  \label{def:physics:phyx-mini-0910:phyxminiproblems-problemphyxmini0910-electricfieldstrengthdimension}
  \lean{PhyXMiniProblems.ProblemPhyXMini0910.electricFieldStrengthDimension}
  The physical dimension M L T⁻² C⁻¹ of electric-field strength
\end{definition}
\begin{proof}
  This is the declared model definition of electric Field Strength Dimension.
\end{proof}
\begin{definition}[force Dimension]
  \label{def:physics:phyx-mini-0910:phyxminiproblems-problemphyxmini0910-forcedimension}
  \lean{PhyXMiniProblems.ProblemPhyXMini0910.forceDimension}
  The physical dimension M L T⁻² of force
\end{definition}
\begin{proof}
  This is the declared model definition of force Dimension.
\end{proof}
\begin{definition}[electric Dipole Moment Dimension]
  \label{def:physics:phyx-mini-0910:phyxminiproblems-problemphyxmini0910-electricdipolemomentdimension}
  \lean{PhyXMiniProblems.ProblemPhyXMini0910.electricDipoleMomentDimension}
  The physical dimension C L of electric-dipole moment
\end{definition}
\begin{proof}
  This is the declared model definition of electric Dipole Moment Dimension.
\end{proof}
\begin{definition}[torque Dimension]
  \label{def:physics:phyx-mini-0910:phyxminiproblems-problemphyxmini0910-torquedimension}
  \lean{PhyXMiniProblems.ProblemPhyXMini0910.torqueDimension}
  The physical dimension M L² T⁻² of torque
\end{definition}
\begin{proof}
  This is the declared model definition of torque Dimension.
\end{proof}
\begin{definition}[moment Of Inertia Dimension]
  \label{def:physics:phyx-mini-0910:phyxminiproblems-problemphyxmini0910-momentofinertiadimension}
  \lean{PhyXMiniProblems.ProblemPhyXMini0910.momentOfInertiaDimension}
  The physical dimension M L² of a scalar moment of inertia
\end{definition}
\begin{proof}
  This is the declared model definition of moment Of Inertia Dimension.
\end{proof}
\begin{definition}[angular Acceleration Dimension]
  \label{def:physics:phyx-mini-0910:phyxminiproblems-problemphyxmini0910-angularaccelerationdimension}
  \lean{PhyXMiniProblems.ProblemPhyXMini0910.angularAccelerationDimension}
  Angular acceleration has dimension T⁻²; radians are dimensionless
\end{definition}
\begin{proof}
  This is the declared model definition of angular Acceleration Dimension.
\end{proof}
\begin{definition}[Mass Quantity]
  \label{def:physics:phyx-mini-0910:phyxminiproblems-problemphyxmini0910-massquantity}
  \lean{PhyXMiniProblems.ProblemPhyXMini0910.MassQuantity}
  A nonnegative, unit-independent physical mass
\end{definition}
\begin{proof}
  This is the declared model definition of Mass Quantity.
\end{proof}
\begin{definition}[Length Quantity]
  \label{def:physics:phyx-mini-0910:phyxminiproblems-problemphyxmini0910-lengthquantity}
  \lean{PhyXMiniProblems.ProblemPhyXMini0910.LengthQuantity}
  A nonnegative, unit-independent physical length
\end{definition}
\begin{proof}
  This is the declared model definition of Length Quantity.
\end{proof}
\begin{definition}[Charge Magnitude Quantity]
  \label{def:physics:phyx-mini-0910:phyxminiproblems-problemphyxmini0910-chargemagnitudequantity}
  \lean{PhyXMiniProblems.ProblemPhyXMini0910.ChargeMagnitudeQuantity}
  A nonnegative, unit-independent electric-charge magnitude
\end{definition}
\begin{proof}
  This is the declared model definition of Charge Magnitude Quantity.
\end{proof}
\begin{definition}[Electric Field Strength Magnitude]
  \label{def:physics:phyx-mini-0910:phyxminiproblems-problemphyxmini0910-electricfieldstrengthmagnitude}
  \lean{PhyXMiniProblems.ProblemPhyXMini0910.ElectricFieldStrengthMagnitude}
  \uses{def:physics:phyx-mini-0910:phyxminiproblems-problemphyxmini0910-electricfieldstrengthdimension}
  A nonnegative, unit-independent electric-field-strength magnitude
\end{definition}
\begin{proof}
  This is the declared model definition of Electric Field Strength Magnitude.
\end{proof}
\begin{definition}[Force Magnitude Quantity]
  \label{def:physics:phyx-mini-0910:phyxminiproblems-problemphyxmini0910-forcemagnitudequantity}
  \lean{PhyXMiniProblems.ProblemPhyXMini0910.ForceMagnitudeQuantity}
  \uses{def:physics:phyx-mini-0910:phyxminiproblems-problemphyxmini0910-forcedimension}
  A nonnegative electric-force magnitude
\end{definition}
\begin{proof}
  This is the declared model definition of Force Magnitude Quantity.
\end{proof}
\begin{definition}[Electric Dipole Moment Magnitude]
  \label{def:physics:phyx-mini-0910:phyxminiproblems-problemphyxmini0910-electricdipolemomentmagnitude}
  \lean{PhyXMiniProblems.ProblemPhyXMini0910.ElectricDipoleMomentMagnitude}
  \uses{def:physics:phyx-mini-0910:phyxminiproblems-problemphyxmini0910-electricdipolemomentdimension}
  A nonnegative electric-dipole-moment magnitude
\end{definition}
\begin{proof}
  This is the declared model definition of Electric Dipole Moment Magnitude.
\end{proof}
\begin{definition}[Torque Magnitude Quantity]
  \label{def:physics:phyx-mini-0910:phyxminiproblems-problemphyxmini0910-torquemagnitudequantity}
  \lean{PhyXMiniProblems.ProblemPhyXMini0910.TorqueMagnitudeQuantity}
  \uses{def:physics:phyx-mini-0910:phyxminiproblems-problemphyxmini0910-torquedimension}
  A nonnegative torque magnitude about the midpoint of the rod
\end{definition}
\begin{proof}
  This is the declared model definition of Torque Magnitude Quantity.
\end{proof}
\begin{definition}[Moment Of Inertia Quantity]
  \label{def:physics:phyx-mini-0910:phyxminiproblems-problemphyxmini0910-momentofinertiaquantity}
  \lean{PhyXMiniProblems.ProblemPhyXMini0910.MomentOfInertiaQuantity}
  \uses{def:physics:phyx-mini-0910:phyxminiproblems-problemphyxmini0910-momentofinertiadimension}
  A nonnegative scalar moment of inertia about the midpoint of the rod
\end{definition}
\begin{proof}
  This is the declared model definition of Moment Of Inertia Quantity.
\end{proof}
\begin{definition}[Angular Acceleration Magnitude]
  \label{def:physics:phyx-mini-0910:phyxminiproblems-problemphyxmini0910-angularaccelerationmagnitude}
  \lean{PhyXMiniProblems.ProblemPhyXMini0910.AngularAccelerationMagnitude}
  \uses{def:physics:phyx-mini-0910:phyxminiproblems-problemphyxmini0910-angularaccelerationdimension}
  A nonnegative angular-acceleration magnitude
\end{definition}
\begin{proof}
  This is the declared model definition of Angular Acceleration Magnitude.
\end{proof}
\begin{definition}[nonnegative SIReadout]
  \label{def:physics:phyx-mini-0910:phyxminiproblems-problemphyxmini0910-nonnegativesireadout}
  \lean{PhyXMiniProblems.ProblemPhyXMini0910.nonnegativeSIReadout}
  Read a nonnegative dimensionful quantity in coherent SI base units
\end{definition}
\begin{proof}
  This is the declared model definition of nonnegative SIReadout.
\end{proof}
\begin{definition}[mass In Kilograms]
  \label{def:physics:phyx-mini-0910:phyxminiproblems-problemphyxmini0910-massinkilograms}
  \lean{PhyXMiniProblems.ProblemPhyXMini0910.massInKilograms}
  \uses{def:physics:phyx-mini-0910:phyxminiproblems-problemphyxmini0910-massquantity, def:physics:phyx-mini-0910:phyxminiproblems-problemphyxmini0910-nonnegativesireadout}
  Read a mass in kilograms
\end{definition}
\begin{proof}
  This is the declared model definition of mass In Kilograms.
\end{proof}
\begin{definition}[mass In Grams]
  \label{def:physics:phyx-mini-0910:phyxminiproblems-problemphyxmini0910-massingrams}
  \lean{PhyXMiniProblems.ProblemPhyXMini0910.massInGrams}
  \uses{def:physics:phyx-mini-0910:phyxminiproblems-problemphyxmini0910-massquantity, def:physics:phyx-mini-0910:phyxminiproblems-problemphyxmini0910-massinkilograms}
  Read a mass in grams, the unit printed beside each ball
\end{definition}
\begin{proof}
  This is the declared model definition of mass In Grams.
\end{proof}
\begin{definition}[length In Meters]
  \label{def:physics:phyx-mini-0910:phyxminiproblems-problemphyxmini0910-lengthinmeters}
  \lean{PhyXMiniProblems.ProblemPhyXMini0910.lengthInMeters}
  \uses{def:physics:phyx-mini-0910:phyxminiproblems-problemphyxmini0910-lengthquantity, def:physics:phyx-mini-0910:phyxminiproblems-problemphyxmini0910-nonnegativesireadout}
  Read a length in metres
\end{definition}
\begin{proof}
  This is the declared model definition of length In Meters.
\end{proof}
\begin{definition}[length In Centimeters]
  \label{def:physics:phyx-mini-0910:phyxminiproblems-problemphyxmini0910-lengthincentimeters}
  \lean{PhyXMiniProblems.ProblemPhyXMini0910.lengthInCentimeters}
  \uses{def:physics:phyx-mini-0910:phyxminiproblems-problemphyxmini0910-lengthquantity, def:physics:phyx-mini-0910:phyxminiproblems-problemphyxmini0910-lengthinmeters}
  Read a length in centimetres, the unit printed beside s
\end{definition}
\begin{proof}
  This is the declared model definition of length In Centimeters.
\end{proof}
\begin{definition}[charge Magnitude In Coulombs]
  \label{def:physics:phyx-mini-0910:phyxminiproblems-problemphyxmini0910-chargemagnitudeincoulombs}
  \lean{PhyXMiniProblems.ProblemPhyXMini0910.chargeMagnitudeInCoulombs}
  \uses{def:physics:phyx-mini-0910:phyxminiproblems-problemphyxmini0910-chargemagnitudequantity, def:physics:phyx-mini-0910:phyxminiproblems-problemphyxmini0910-nonnegativesireadout}
  Read a charge magnitude in coulombs
\end{definition}
\begin{proof}
  This is the declared model definition of charge Magnitude In Coulombs.
\end{proof}
\begin{definition}[charge Magnitude In Nanocoulombs]
  \label{def:physics:phyx-mini-0910:phyxminiproblems-problemphyxmini0910-chargemagnitudeinnanocoulombs}
  \lean{PhyXMiniProblems.ProblemPhyXMini0910.chargeMagnitudeInNanocoulombs}
  \uses{def:physics:phyx-mini-0910:phyxminiproblems-problemphyxmini0910-chargemagnitudequantity, def:physics:phyx-mini-0910:phyxminiproblems-problemphyxmini0910-chargemagnitudeincoulombs}
  Read a charge magnitude in nanocoulombs
\end{definition}
\begin{proof}
  This is the declared model definition of charge Magnitude In Nanocoulombs.
\end{proof}
\begin{definition}[electric Field Strength In Newtons Per Coulomb]
  \label{def:physics:phyx-mini-0910:phyxminiproblems-problemphyxmini0910-electricfieldstrengthinnewtonspercoulomb}
  \lean{PhyXMiniProblems.ProblemPhyXMini0910.electricFieldStrengthInNewtonsPerCoulomb}
  \uses{def:physics:phyx-mini-0910:phyxminiproblems-problemphyxmini0910-electricfieldstrengthmagnitude, def:physics:phyx-mini-0910:phyxminiproblems-problemphyxmini0910-nonnegativesireadout}
  Read electric-field strength in newtons per coulomb
\end{definition}
\begin{proof}
  This is the declared model definition of electric Field Strength In Newtons Per Coulomb.
\end{proof}
\begin{definition}[force Magnitude In Newtons]
  \label{def:physics:phyx-mini-0910:phyxminiproblems-problemphyxmini0910-forcemagnitudeinnewtons}
  \lean{PhyXMiniProblems.ProblemPhyXMini0910.forceMagnitudeInNewtons}
  \uses{def:physics:phyx-mini-0910:phyxminiproblems-problemphyxmini0910-forcemagnitudequantity, def:physics:phyx-mini-0910:phyxminiproblems-problemphyxmini0910-nonnegativesireadout}
  Read an electric-force magnitude in newtons
\end{definition}
\begin{proof}
  This is the declared model definition of force Magnitude In Newtons.
\end{proof}
\begin{definition}[dipole Moment In Coulomb Meters]
  \label{def:physics:phyx-mini-0910:phyxminiproblems-problemphyxmini0910-dipolemomentincoulombmeters}
  \lean{PhyXMiniProblems.ProblemPhyXMini0910.dipoleMomentInCoulombMeters}
  \uses{def:physics:phyx-mini-0910:phyxminiproblems-problemphyxmini0910-electricdipolemomentmagnitude, def:physics:phyx-mini-0910:phyxminiproblems-problemphyxmini0910-nonnegativesireadout}
  Read an electric-dipole moment magnitude in coulomb-metres
\end{definition}
\begin{proof}
  This is the declared model definition of dipole Moment In Coulomb Meters.
\end{proof}
\begin{definition}[torque In Newton Meters]
  \label{def:physics:phyx-mini-0910:phyxminiproblems-problemphyxmini0910-torqueinnewtonmeters}
  \lean{PhyXMiniProblems.ProblemPhyXMini0910.torqueInNewtonMeters}
  \uses{def:physics:phyx-mini-0910:phyxminiproblems-problemphyxmini0910-torquemagnitudequantity, def:physics:phyx-mini-0910:phyxminiproblems-problemphyxmini0910-nonnegativesireadout}
  Read a torque magnitude in newton-metres
\end{definition}
\begin{proof}
  This is the declared model definition of torque In Newton Meters.
\end{proof}
\begin{definition}[moment Of Inertia In Kilogram Meters Squared]
  \label{def:physics:phyx-mini-0910:phyxminiproblems-problemphyxmini0910-momentofinertiainkilogrammeterssquared}
  \lean{PhyXMiniProblems.ProblemPhyXMini0910.momentOfInertiaInKilogramMetersSquared}
  \uses{def:physics:phyx-mini-0910:phyxminiproblems-problemphyxmini0910-momentofinertiaquantity, def:physics:phyx-mini-0910:phyxminiproblems-problemphyxmini0910-nonnegativesireadout}
  Read a moment of inertia in kilogram-metres squared
\end{definition}
\begin{proof}
  This is the declared model definition of moment Of Inertia In Kilogram Meters Squared.
\end{proof}
\begin{definition}[angular Acceleration In Radians Per Second Squared]
  \label{def:physics:phyx-mini-0910:phyxminiproblems-problemphyxmini0910-angularaccelerationinradianspersecondsquared}
  \lean{PhyXMiniProblems.ProblemPhyXMini0910.angularAccelerationInRadiansPerSecondSquared}
  \uses{def:physics:phyx-mini-0910:phyxminiproblems-problemphyxmini0910-angularaccelerationmagnitude, def:physics:phyx-mini-0910:phyxminiproblems-problemphyxmini0910-nonnegativesireadout}
  Read angular acceleration in radians per second squared
\end{definition}
\begin{proof}
  This is the declared model definition of angular Acceleration In Radians Per Second Squared.
\end{proof}
\begin{definition}[degrees To Angle]
  \label{def:physics:phyx-mini-0910:phyxminiproblems-problemphyxmini0910-degreestoangle}
  \lean{PhyXMiniProblems.ProblemPhyXMini0910.degreesToAngle}
  Convert the degree measure used in the image to Mathlib's angle type
\end{definition}
\begin{proof}
  This is the declared model definition of degrees To Angle.
\end{proof}
\begin{definition}[Ball Label]
  \label{def:physics:phyx-mini-0910:phyxminiproblems-problemphyxmini0910-balllabel}
  \lean{PhyXMiniProblems.ProblemPhyXMini0910.BallLabel}
  The two charged balls, named by the signs printed in the image
\end{definition}
\begin{proof}
  This is the declared model definition of Ball Label.
\end{proof}
\begin{definition}[Charge Sign]
  \label{def:physics:phyx-mini-0910:phyxminiproblems-problemphyxmini0910-chargesign}
  \lean{PhyXMiniProblems.ProblemPhyXMini0910.ChargeSign}
  The sign carried by a ball's electric charge
\end{definition}
\begin{proof}
  This is the declared model definition of Charge Sign.
\end{proof}
\begin{definition}[Ball Placement]
  \label{def:physics:phyx-mini-0910:phyxminiproblems-problemphyxmini0910-ballplacement}
  \lean{PhyXMiniProblems.ProblemPhyXMini0910.BallPlacement}
  Coarse ball locations in the supplied planar drawing
\end{definition}
\begin{proof}
  This is the declared model definition of Ball Placement.
\end{proof}
\begin{definition}[Planar Direction]
  \label{def:physics:phyx-mini-0910:phyxminiproblems-problemphyxmini0910-planardirection}
  \lean{PhyXMiniProblems.ProblemPhyXMini0910.PlanarDirection}
  Qualitative directions of arrows in the supplied planar drawing
\end{definition}
\begin{proof}
  This is the declared model definition of Planar Direction.
\end{proof}
\begin{definition}[External Field Model]
  \label{def:physics:phyx-mini-0910:phyxminiproblems-problemphyxmini0910-externalfieldmodel}
  \lean{PhyXMiniProblems.ProblemPhyXMini0910.ExternalFieldModel}
  The idealized external electric-field models relevant to the problem
\end{definition}
\begin{proof}
  This is the declared model definition of External Field Model.
\end{proof}
\begin{definition}[Rod Mass Model]
  \label{def:physics:phyx-mini-0910:phyxminiproblems-problemphyxmini0910-rodmassmodel}
  \lean{PhyXMiniProblems.ProblemPhyXMini0910.RodMassModel}
  The mass model for the insulating connecting rod
\end{definition}
\begin{proof}
  This is the declared model definition of Rod Mass Model.
\end{proof}
\begin{definition}[Endpoint Mass Model]
  \label{def:physics:phyx-mini-0910:phyxminiproblems-problemphyxmini0910-endpointmassmodel}
  \lean{PhyXMiniProblems.ProblemPhyXMini0910.EndpointMassModel}
  The model used for rotational inertia when no ball radii are supplied
\end{definition}
\begin{proof}
  This is the declared model definition of Endpoint Mass Model.
\end{proof}
\begin{definition}[Release Protocol]
  \label{def:physics:phyx-mini-0910:phyxminiproblems-problemphyxmini0910-releaseprotocol}
  \lean{PhyXMiniProblems.ProblemPhyXMini0910.ReleaseProtocol}
  How the rod-and-ball assembly begins its motion
\end{definition}
\begin{proof}
  This is the declared model definition of Release Protocol.
\end{proof}
\begin{definition}[force Direction In Rightward Field]
  \label{def:physics:phyx-mini-0910:phyxminiproblems-problemphyxmini0910-forcedirectioninrightwardfield}
  \lean{PhyXMiniProblems.ProblemPhyXMini0910.forceDirectionInRightwardField}
  \uses{def:physics:phyx-mini-0910:phyxminiproblems-problemphyxmini0910-chargesign, def:physics:phyx-mini-0910:phyxminiproblems-problemphyxmini0910-planardirection}
  Direction of the force on a signed charge in a rightward electric field
\end{definition}
\begin{proof}
  This is the declared model definition of force Direction In Rightward Field.
\end{proof}
\begin{definition}[Electric Dipole Figure]
  \label{def:physics:phyx-mini-0910:phyxminiproblems-problemphyxmini0910-electricdipolefigure}
  \lean{PhyXMiniProblems.ProblemPhyXMini0910.ElectricDipoleFigure}
  \uses{def:physics:phyx-mini-0910:phyxminiproblems-problemphyxmini0910-balllabel, def:physics:phyx-mini-0910:phyxminiproblems-problemphyxmini0910-chargesign, def:physics:phyx-mini-0910:phyxminiproblems-problemphyxmini0910-ballplacement, def:physics:phyx-mini-0910:phyxminiproblems-problemphyxmini0910-planardirection}
  Literal visual information from image 910.png. Numerical entries are the values printed in their displayed units, while the physical quantities remain separate fields of the setup below
\end{definition}
\begin{proof}
  This is the declared model definition of Electric Dipole Figure.
\end{proof}
\begin{definition}[Electric Dipole Release Setup]
  \label{def:physics:phyx-mini-0910:phyxminiproblems-problemphyxmini0910-electricdipolereleasesetup}
  \lean{PhyXMiniProblems.ProblemPhyXMini0910.ElectricDipoleReleaseSetup}
  \uses{def:physics:phyx-mini-0910:phyxminiproblems-problemphyxmini0910-massquantity, def:physics:phyx-mini-0910:phyxminiproblems-problemphyxmini0910-lengthquantity, def:physics:phyx-mini-0910:phyxminiproblems-problemphyxmini0910-chargemagnitudequantity, def:physics:phyx-mini-0910:phyxminiproblems-problemphyxmini0910-electricfieldstrengthmagnitude, def:physics:phyx-mini-0910:phyxminiproblems-problemphyxmini0910-forcemagnitudequantity, def:physics:phyx-mini-0910:phyxminiproblems-problemphyxmini0910-electricdipolemomentmagnitude, def:physics:phyx-mini-0910:phyxminiproblems-problemphyxmini0910-torquemagnitudequantity, def:physics:phyx-mini-0910:phyxminiproblems-problemphyxmini0910-momentofinertiaquantity, def:physics:phyx-mini-0910:phyxminiproblems-problemphyxmini0910-angularaccelerationmagnitude, def:physics:phyx-mini-0910:phyxminiproblems-problemphyxmini0910-balllabel, def:physics:phyx-mini-0910:phyxminiproblems-problemphyxmini0910-chargesign, def:physics:phyx-mini-0910:phyxminiproblems-problemphyxmini0910-planardirection, def:physics:phyx-mini-0910:phyxminiproblems-problemphyxmini0910-externalfieldmodel, def:physics:phyx-mini-0910:phyxminiproblems-problemphyxmini0910-rodmassmodel, def:physics:phyx-mini-0910:phyxminiproblems-problemphyxmini0910-endpointmassmodel, def:physics:phyx-mini-0910:phyxminiproblems-problemphyxmini0910-releaseprotocol, def:physics:phyx-mini-0910:phyxminiproblems-problemphyxmini0910-electricdipolefigure}
  Independent physical quantities for the release configuration. In particular, the initial angular acceleration is not defined from any answer choice; it is related to torque and inertia only by the law structure below
\end{definition}
\begin{proof}
  This is the declared model definition of Electric Dipole Release Setup.
\end{proof}
\begin{definition}[Matches Primary Figure]
  \label{def:physics:phyx-mini-0910:phyxminiproblems-problemphyxmini0910-matchesprimaryfigure}
  \lean{PhyXMiniProblems.ProblemPhyXMini0910.MatchesPrimaryFigure}
  \uses{def:physics:phyx-mini-0910:phyxminiproblems-problemphyxmini0910-electricdipolereleasesetup}
  The positions, labels, and arrows read directly from the primary raster. This record does not constrain the requested angular acceleration
\end{definition}
\begin{proof}
  This is the declared model definition of Matches Primary Figure.
\end{proof}
\begin{definition}[Matches Problem Data]
  \label{def:physics:phyx-mini-0910:phyxminiproblems-problemphyxmini0910-matchesproblemdata}
  \lean{PhyXMiniProblems.ProblemPhyXMini0910.MatchesProblemData}
  \uses{def:physics:phyx-mini-0910:phyxminiproblems-problemphyxmini0910-massingrams, def:physics:phyx-mini-0910:phyxminiproblems-problemphyxmini0910-lengthincentimeters, def:physics:phyx-mini-0910:phyxminiproblems-problemphyxmini0910-chargemagnitudeinnanocoulombs, def:physics:phyx-mini-0910:phyxminiproblems-problemphyxmini0910-electricfieldstrengthinnewtonspercoulomb, def:physics:phyx-mini-0910:phyxminiproblems-problemphyxmini0910-degreestoangle, def:physics:phyx-mini-0910:phyxminiproblems-problemphyxmini0910-electricdipolereleasesetup}
  Physical numerical readouts and idealizations stated by the problem. The point-mass model places each labeled ball mass at its center because the problem supplies no ball radius and declares the rod mass negligible
\end{definition}
\begin{proof}
  This is the declared model definition of Matches Problem Data.
\end{proof}
\begin{definition}[Satisfies Electric Dipole Rotational Laws]
  \label{def:physics:phyx-mini-0910:phyxminiproblems-problemphyxmini0910-satisfieselectricdipolerotationallaws}
  \lean{PhyXMiniProblems.ProblemPhyXMini0910.SatisfiesElectricDipoleRotationalLaws}
  \uses{def:physics:phyx-mini-0910:phyxminiproblems-problemphyxmini0910-massinkilograms, def:physics:phyx-mini-0910:phyxminiproblems-problemphyxmini0910-lengthinmeters, def:physics:phyx-mini-0910:phyxminiproblems-problemphyxmini0910-chargemagnitudeincoulombs, def:physics:phyx-mini-0910:phyxminiproblems-problemphyxmini0910-electricfieldstrengthinnewtonspercoulomb, def:physics:phyx-mini-0910:phyxminiproblems-problemphyxmini0910-forcemagnitudeinnewtons, def:physics:phyx-mini-0910:phyxminiproblems-problemphyxmini0910-dipolemomentincoulombmeters, def:physics:phyx-mini-0910:phyxminiproblems-problemphyxmini0910-torqueinnewtonmeters, def:physics:phyx-mini-0910:phyxminiproblems-problemphyxmini0910-momentofinertiainkilogrammeterssquared, def:physics:phyx-mini-0910:phyxminiproblems-problemphyxmini0910-angularaccelerationinradianspersecondsquared, def:physics:phyx-mini-0910:phyxminiproblems-problemphyxmini0910-forcedirectioninrightwardfield, def:physics:phyx-mini-0910:phyxminiproblems-problemphyxmini0910-electricdipolereleasesetup}
  Electrostatic and rotational laws for the idealized release configuration. They relate independently stored observables and contain neither the target value 5 nor any answer-choice label
\end{definition}
\begin{proof}
  This is the declared model definition of Satisfies Electric Dipole Rotational Laws.
\end{proof}
\begin{lemma}[initial Electric Torque In Newton Meters eq one Millionth]
  \label{lem:physics:phyx-mini-0910:phyxminiproblems-problemphyxmini0910-initialelectrictorqueinnewtonmeters-eq-onemillionth}
  \lean{PhyXMiniProblems.ProblemPhyXMini0910.initialElectricTorqueInNewtonMeters_eq_oneMillionth}
  \uses{def:physics:phyx-mini-0910:phyxminiproblems-problemphyxmini0910-torqueinnewtonmeters, def:physics:phyx-mini-0910:phyxminiproblems-problemphyxmini0910-electricdipolereleasesetup, def:physics:phyx-mini-0910:phyxminiproblems-problemphyxmini0910-matchesproblemdata, def:physics:phyx-mini-0910:phyxminiproblems-problemphyxmini0910-satisfieselectricdipolerotationallaws}
  The applied electric couple initially has magnitude 10⁻⁶ N m
\end{lemma}
\begin{proof}
  The conclusion follows from the governing relations and figure readouts named in the statement of initial Electric Torque In Newton Meters eq one Millionth.
\end{proof}
\begin{lemma}[moment Of Inertia In Kilogram Meters Squared eq two Tenths Millionth]
  \label{lem:physics:phyx-mini-0910:phyxminiproblems-problemphyxmini0910-momentofinertiainkilogrammeterssquared-eq-twotenthsmillionth}
  \lean{PhyXMiniProblems.ProblemPhyXMini0910.momentOfInertiaInKilogramMetersSquared_eq_twoTenthsMillionth}
  \uses{def:physics:phyx-mini-0910:phyxminiproblems-problemphyxmini0910-momentofinertiainkilogrammeterssquared, def:physics:phyx-mini-0910:phyxminiproblems-problemphyxmini0910-electricdipolereleasesetup, def:physics:phyx-mini-0910:phyxminiproblems-problemphyxmini0910-matchesproblemdata, def:physics:phyx-mini-0910:phyxminiproblems-problemphyxmini0910-satisfieselectricdipolerotationallaws}
  The two 1 g endpoint masses have axial inertia 2 * 10⁻⁷ kg m²
\end{lemma}
\begin{proof}
  The conclusion follows from the governing relations and figure readouts named in the statement of moment Of Inertia In Kilogram Meters Squared eq two Tenths Millionth.
\end{proof}
\begin{definition}[Answer Choice]
  \label{def:physics:phyx-mini-0910:phyxminiproblems-problemphyxmini0910-answerchoice}
  \lean{PhyXMiniProblems.ProblemPhyXMini0910.AnswerChoice}
  Labels of the four angular-acceleration choices in the source
\end{definition}
\begin{proof}
  This is the declared model definition of Answer Choice.
\end{proof}
\begin{definition}[displayed Angular Acceleration In Radians Per Second Squared]
  \label{def:physics:phyx-mini-0910:phyxminiproblems-problemphyxmini0910-displayedangularaccelerationinradianspersecondsquared}
  \lean{PhyXMiniProblems.ProblemPhyXMini0910.displayedAngularAccelerationInRadiansPerSecondSquared}
  \uses{def:physics:phyx-mini-0910:phyxminiproblems-problemphyxmini0910-answerchoice}
  Radian-per-second-squared readout printed beside each choice
\end{definition}
\begin{proof}
  This is the declared model definition of displayed Angular Acceleration In Radians Per Second Squared.
\end{proof}
\begin{definition}[recorded Dataset Answer]
  \label{def:physics:phyx-mini-0910:phyxminiproblems-problemphyxmini0910-recordeddatasetanswer}
  \lean{PhyXMiniProblems.ProblemPhyXMini0910.recordedDatasetAnswer}
  \uses{def:physics:phyx-mini-0910:phyxminiproblems-problemphyxmini0910-answerchoice}
  The answer label recorded by the source dataset; it is not a premise
\end{definition}
\begin{proof}
  This is the declared model definition of recorded Dataset Answer.
\end{proof}
\begin{definition}[Answer Matches Initial Angular Acceleration]
  \label{def:physics:phyx-mini-0910:phyxminiproblems-problemphyxmini0910-answermatchesinitialangularacceleration}
  \lean{PhyXMiniProblems.ProblemPhyXMini0910.AnswerMatchesInitialAngularAcceleration}
  \uses{def:physics:phyx-mini-0910:phyxminiproblems-problemphyxmini0910-angularaccelerationinradianspersecondsquared, def:physics:phyx-mini-0910:phyxminiproblems-problemphyxmini0910-electricdipolereleasesetup, def:physics:phyx-mini-0910:phyxminiproblems-problemphyxmini0910-answerchoice, def:physics:phyx-mini-0910:phyxminiproblems-problemphyxmini0910-displayedangularaccelerationinradianspersecondsquared}
  A displayed choice agrees with the independently modeled acceleration
\end{definition}
\begin{proof}
  This is the declared model definition of Answer Matches Initial Angular Acceleration.
\end{proof}
% --- Archon declaration topology end ---
% --- Archon physics formalization source end ---
